\chapter{Odpowiedniość Galois}

\section{Niezależność zanurzeń i twierdzenie Artina}

\begin{lemat}[Liniowa niezależność homomorfizmów]
Różne homomorfizmy ciał $\sigma_1,\dots,\sigma_m:L\to M$ są liniowo niezależne nad $M$ jako funkcje: z $\sum_{j=1}^m c_j\sigma_j(x)=0$ dla każdego $x\in L$ wynika $c_1=\cdots=c_m=0$.
\end{lemat}
\begin{proof}
Przypuśćmy, że istnieje niezerowa relacja, i wybierzmy taką z najmniejszą liczbą niezerowych współczynników. Jeden składnik nie wystarcza, co widać dla $x=1$. Numerując składniki, możemy założyć $c_1,c_m\ne0$. Ponieważ $\sigma_1\ne\sigma_m$, istnieje $a\in L$ takie, że $\sigma_1(a)\ne\sigma_m(a)$. Od relacji obliczonej w $ax$ odejmujemy relację w $x$ pomnożoną przez $\sigma_m(a)$. Dostajemy
\[
\sum_{j=1}^{m-1}c_j\bigl(\sigma_j(a)-\sigma_m(a)\bigr)\sigma_j(x)=0.
\]
Współczynnik pierwszego składnika jest niezerowy. To relacja z mniejszą liczbą składników, wbrew wyborowi.
\end{proof}

\begin{twierdzenie}[Artina: stopień nad ciałem stałym]
Niech $H$ będzie skończoną grupą automorfizmów ciała $L$, a
\[
F=L^H=\{a\in L:\sigma(a)=a\text{ dla każdego }\sigma\in H\}.
\]
Wtedy $[L:F]=|H|$ oraz $\Aut_F(L)=H$.
\end{twierdzenie}
\begin{proof}
Niech $m=|H|$. Najpierw wykażemy $[L:F]\le m$, nie zakładając z góry skończoności tego stopnia. Dla dowolnych $x_1,\dots,x_{m+1}\in L$ układ $m$ równań
\[
\sum_{j=1}^{m+1}c_j\sigma(x_j)=0\qquad(\sigma\in H)
\]
ma niezerowe rozwiązanie $(c_j)$ nad $L$. Wybierzmy rozwiązanie o najmniejszej liczbie niezerowych współrzędnych i unormujmy jedną z nich do $c_1=1$, w razie potrzeby zmieniając numerację. Dla $\tau\in H$ wektor $(\tau(c_j))$ też jest rozwiązaniem: zastosowanie $\tau$ do układu jedynie permutuje wiersze $\tau\sigma$. Różnica $(\tau(c_j)-c_j)$ ma nośnik zawarty w nośniku $(c_j)$ i zerową pierwszą współrzędną. Minimalność wymusza, że jest zerowa. Wszystkie $c_j$ należą zatem do $F$. Wiersz odpowiadający identyczności daje niezerową zależność $F$-liniową między $x_j$. Stąd $[L:F]\le m$.

Oznaczmy $d=[L:F]$. Przestrzeń $F$-liniowych odwzorowań $L\to L$ ma wymiar $d$ nad $L$: odwzorowanie jest wyznaczone dowolnymi obrazami $d$ elementów bazy. Lemat 3.1 daje $m\le d$, więc $d=m$. Każde $d+1$ różnych $F$-automorfizmów przeczyłoby temu samemu lematowi. Ponieważ $H\subseteq\Aut_F(L)$ i $|H|=d$, zachodzi $\Aut_F(L)=H$.
\end{proof}

\section{Twierdzenie główne z pełnym dowodem}

\begin{twierdzenie}[Podstawowe twierdzenie teorii Galois]
Niech $L/K$ będzie skończonym rozszerzeniem Galois w charakterystyce zero, $G=\Gal(L/K)$. Odwzorowania
\[
E\longmapsto\Gal(L/E),\qquad H\longmapsto L^H
\]
są wzajemnie odwrotnymi bijekcjami między ciałami pośrednimi a podgrupami $G$, odwracają inkluzje i spełniają
\[
[L:E]=|\Gal(L/E)|,\qquad [E:K]=|G:\Gal(L/E)|.
\]
Jeśli $H=\Gal(L/E)$, to $E/K$ jest Galois wtedy i tylko wtedy, gdy $H\triangleleft G$. W tym przypadku ograniczenie automorfizmów daje $\Gal(E/K)\simeq G/H$.
\end{twierdzenie}
\begin{proof}
\emph{Ciała odzyskujemy z podgrup.} Ciało $L$ jest ciałem rozkładu pewnego $f\in K[T]$, więc jest też ciałem rozkładu $f$ nad każdym $E$. Stąd $L/E$ jest Galois i $|\Gal(L/E)|=[L:E]$. Przy $H=\Gal(L/E)$ mamy $E\subseteq L^H$, a twierdzenie Artina daje $[L:L^H]=|H|=[L:E]$. Twierdzenie o wieży wymusza $E=L^H$. W szczególności $L^G=K$.

\emph{Podgrupy odzyskujemy z ciał.} Dla $H\le G$ twierdzenie Artina daje bezpośrednio $\Gal(L/L^H)=H$. Są to więc bijekcje odwrotne. Większa podgrupa ustala mniej elementów, a większe ciało pozostawia mniej automorfizmów; stąd odwrócenie inkluzji. Wzór na $[E:K]$ wynika z twierdzenia o wieży oraz $[L:K]=|G|$.

\emph{Normalność.} Dla $\sigma\in G$ bezpośrednio z definicji wynika
\[
\sigma(L^H)=L^{\sigma H\sigma^{-1}}.
\]
Istotnie, $h(a)=a$ jest równoważne $(\sigma h\sigma^{-1})(\sigma(a))=\sigma(a)$. Jeśli $H$ jest normalna, to każde $\sigma$ zachowuje $E=L^H$. Dowolne $K$-zanurzenie $E\to\overline K$ przedłuża się do $L$; normalność $L/K$ sprawia, że przedłużenie należy do $G$. Każde takie zanurzenie ma więc obraz $E$. Lemat 2.3 daje normalność $E/K$. Odwrotnie, gdy $E/K$ jest normalne, $\sigma(E)=E$ dla każdego $\sigma\in G$. Wzór powyżej i bijekcja podgrup z ciałami dają $\sigma H\sigma^{-1}=H$.

\emph{Grupa ilorazowa.} Przy normalności ograniczenie $G\to\Gal(E/K)$ jest homomorfizmem z jądrem $H$. Jest surjektywne, ponieważ każde zanurzenie $E$ przedłuża się do $L$. Twierdzenie o homomorfizmie grup daje $\Gal(E/K)\simeq G/H$.
\end{proof}

\section{Jak czytać odpowiedniość na przykładach}

Dla $L=\Q(\sqrt2,\sqrt3)$ niech $s$ zmienia znak $\sqrt2$, a $t$ znak $\sqrt3$. Wszystkie podgrupy i odpowiadające im ciała to
\[
\begin{array}{c|ccccc}
H & \{1\} & \langle s\rangle & \langle t\rangle & \langle st\rangle & G\\ \hline
L^H & L & \Q(\sqrt3) & \Q(\sqrt2) & \Q(\sqrt6) & \Q.
\end{array}
\]
Na przykład $\sqrt6$ jest ustalane przez $st$, a jego ciało ma stopień $2$; wzór na stopnie dowodzi, że jest to całe ciało stałe, nie tylko jego podciało. Wszystkie podgrupy są normalne, więc wszystkie rozszerzenia pośrednie nad $\Q$ są Galois.

Dla $T^3-2$ grupa $S_3$ ma trzy podgrupy rzędu $2$. Każda ustala jedno zero i odpowiada generowanemu przez nie ciału sześciennemu. Te podgrupy nie są normalne i ciała sześcienne nie są normalne nad $\Q$. Podgrupa $A_3=\langle\sigma\rangle$ jest normalna i ustala $\Q(\omega)$: ciało to ma stopień $2$ i $\sigma(\omega)=\omega$. Tak odróżniamy stopień rozszerzenia od liczby jego automorfizmów, gdy nie ma normalności.

\section{Wyróżnik: szybki test dla grupy wielomianu}

Niech unormowany rozdzielczy $f\in K[T]$ ma zera $\alpha_1,\dots,\alpha_n$ w ciele rozkładu $L$. Połóżmy
\[
\delta=\prod_{i<j}(\alpha_i-\alpha_j),\qquad \Delta(f)=\delta^2.
\]
Zamiana dwóch zer zmienia znak $\delta$; każda permutacja daje zatem $\sigma(\delta)=\operatorname{sgn}(\sigma)\delta$. Kwadrat $\delta^2$ jest ustalany przez całą grupę, więc należy do $K$. Ponieważ $\delta\ne0$ i $\charac K=0$, otrzymujemy
\[
\Gal(f/K)\subseteq A_n\quad\Longleftrightarrow\quad \delta\in K
\quad\Longleftrightarrow\quad \Delta(f)\text{ jest kwadratem w }K.
\]
W ostatnim kroku, jeśli $\Delta=b^2$ dla $b\in K$, to $(\delta-b)(\delta+b)=0$ i $\delta=\pm b$. Warunek rozdzielczości jest istotny: dla $\Delta=0$ ten argument nie wykrywa parzystości permutacji.

\section{Ślad i norma: przygotowanie do całkowania}

Dla skończonego $L/K$ w charakterystyce zero i $d=[L:K]$ definiujemy
\[
\Tr_{L/K}(a)=\sum_{j=1}^d\sigma_j(a),\qquad
\Nm_{L/K}(a)=\prod_{j=1}^d\sigma_j(a),
\]
gdzie $\sigma_j$ są wszystkimi $K$-zanurzeniami $L$ do wspólnego domknięcia algebraicznego. Nie ograniczamy sumy do automorfizmów samego $L$, których w rozszerzeniu nienormalnym może być mniej.

Oba wyrażenia należą do $K$. Aby to uzasadnić, weźmy ciało rozkładu iloczynu wielomianów minimalnych skończonej rodziny generatorów $L$. To skończone rozszerzenie Galois $N/K$ zawiera wszystkie obrazy $L$. Automorfizmy $N/K$ permutują zanurzenia $\sigma_j$, zatem ustalają sumę i iloczyn. Równość $N^{\Gal(N/K)}=K$ dowodzi tezy.

Ślad jest $K$-liniowy, norma jest multiplikatywna i norma elementu niezerowego jest niezerowa. Dla $a\in K$ mamy $\Tr_{L/K}(a)=da$ oraz $\Nm_{L/K}(a)=a^d$. W części o całkach udowodnimy, że ślad komutuje z pochodną, a norma zamienia ślad pochodnej logarytmicznej w pochodną logarytmiczną normy.


